\section{Kontrol Akses dan \textit{Attribute-Based Access Control} (ABAC)}

Kontrol akses merupakan mekanisme keamanan yang menetukan entitas yang berhak melakukan operasi tertentu terhadap suatu sumber daya beserta kondisi yang harus dipenuhi.
Model kontrol akses klasik meliputi \textit{Discretionary Access Control} (DAC), \textit{Mandatory Access Control} (MAC)\label{first:mac}, dan \textit{Role-Based Access Control} (RBAC).
Pada DAC, pemilik sumber daya mengatur hak akses pengguna, umumnya melalui \textit{Access Control List} (ACL)\label{first:acl} yang perlu diperbarui ketika terjadi perubahan.
Sifatnya yang relatif statis membuat DAC kurang sesuai untuk lingkungan IoT yang dinamis \parencite{ragothaman_access_2023,ahsan_comprehensive_2025}.

Berbeda dengan DAC, MAC menentukan akses melalui kebijakan terpusat berbasis label keamanan, seperti klasifikasi objek dan tingkat kewenangan subjek.
Pendekatan ini memberikan kontrol yang ketat dan konsisten, tetapi pengelolaan label yang terpusat dan kaku menyulitkan penerapannya pada lingkungan IoT yang heterogen, berskala besar, dan terus berubah \parencite{ragothaman_access_2023}.

RBAC mengaitkan hak akses dengan peran pengguna sehingga pengelolaannya lebih sederhana dibandingkan pemberian izin per individu.
Namun, pada lingkungan IoT yang sangat heterogen, kebutuhan untuk merepresentasikan berbagai kombinasi konteks dapat memicu \textit{role explosion}.
Selain itu, peran kurang fleksibel untuk merepresentasikan kondisi kontekstual seperti waktu, lokasi, dan status perangkat \parencite{ahsan_comprehensive_2025}.

Sebagai alternatif, \textit{Attribute-Based Access Control} (ABAC) menentukan otorisasi dengan mengevaluasi atribut subjek, objek, operasi, dan lingkungan terhadap kebijakan yang berlaku.
Menurut NIST\label{first:nist} SP 800-162, pendekatan ini memungkinkan kombinasi input keputusan yang lebih luas serta mampu mengakomodasi konsep DAC dan MAC dalam suatu model berbasis atribut \parencite{hu_guide_2014}.

Secara operasional, ABAC terdiri atas tiga komponen utama, yaitu atribut, aturan (\textit{rule}), dan kebijakan \textit{policy}.
Atribut merepresentasikan karakteristik subjek, objek, dan lingkungan; aturan menentukan kondisi pemberian akses berdasarkan atribut tersebut; sedangkan kebijakan merupakan himpunan aturan otorisasi.
Ketika permintaan akses diterima, atribut yang relevan dievaluasi terhadap kebijakan untuk menghasilkan keputusan izin atau penolakan akses \parencite{diaz-rodriguez_abac_2025}.
Alur pengambilan keputusan tersebut diilustrasikan pada Gambar \ref{fig:arsitektur-abac}.

\begin{figure}[H]
    \centering
    \includegraphics[width=0.9\textwidth]{gambar_abac.png}
    \caption{Arsitektur Keputusan Akses ABAC}
    \label{fig:arsitektur-abac}
\end{figure}

Gambar \ref{fig:arsitektur-abac} menunjukkan bahwa keputusan ABAC tidak bergantung langsung pada identitas subjek, tetapi pada kombinasi atribut subjek, objek, dan lingkungan yang dievaluasi terhadap aturan dalam kebijakan.
Struktur ini memungkinkan ABAC merepresentasikan konteks akses yang kompleks, sekaligus menunjukkan bahwa kualitas keputusan akses bergantung pada kualitas aturan yang membentuk kebijakan.

Fleksibilitas tersebut menjadikan ABAC sesuai untuk lingkungan IoT yang heterogen dan dinamis karena keputusan akses dapat mempertimbangkan atribut perangkat, pengguna, dan lingkungan tanpa mendefinisikan peran untuk setiap kemungkinan konteks \parencite{ahsan_comprehensive_2025,ragothaman_access_2023}.
Namun, fleksibilitas ini juga meningkatkan kompleksitas spesifikasi kebijakan, terutama ketika kebijakan dirancang secara manual \parencite{diaz-rodriguez_abac_2025}.